\section{Variational Structure and the KKT Condition}\label{app:KKT_Variational}

This appendix develops the rigorous variational framework for the outermost MOTS, addressing the analytical obstruction identified in Section \ref{sec:Jang}. We shift the perspective from the standard "stability" analysis (based on the principal eigenvalue) to a "constrained maximization" analysis (based on the Karush-Kuhn-Tucker conditions). This approach yields a much stronger structural characterization of the mean curvature jump, sufficient to establish the distributional non-negativity of scalar curvature required for the smoothing argument.

\subsection{The Constrained Maximization Problem}

Let $(\Sigma, \gamma)$ be a closed 2-surface in the initial data set $(M, g, k)$. The outward null expansion is given by $\theta_+ = H_\Sigma + \tr_\Sigma k$.
We consider the problem of maximizing the area of $\Sigma$ subject to the constraint that $\Sigma$ remains a trapped surface ($\theta_+ \le 0$).

\textbf{Multi-Component Convention:} We allow $\Sigma$ to be disconnected, i.e., $\Sigma = \bigcup_i \Sigma_i$. In this case, "Area" refers to the sum of the areas of the connected components, $\mathcal{A}(\Sigma) = \sum_i |\Sigma_i|$, and the constraint $\theta_+ \le 0$ is imposed pointwise on each component. This is crucial for handling multi-black hole configurations where the maximizer may be the union of individual horizons.

Let $\mathcal{S}$ be the space of smooth surfaces homologous to the outer boundary. We define the functional:
\begin{equation}
    \mathcal{A}(\Sigma) = \int_\Sigma dA.
\end{equation}
The constraint is:
\begin{equation}
    \mathcal{C}(\Sigma) := \theta_+(\Sigma) \le 0 \quad \text{pointwise on } \Sigma.
\end{equation}

\subsection{The Stability Operator and Linearization}

Let $\Sigma_t$ be a variation of $\Sigma$ generated by the normal vector field $\phi \nu$. The first variation of area is:
\begin{equation}
    \delta \mathcal{A}[\phi] = \int_\Sigma H_\Sigma \phi \, dA.
\end{equation}
The linearization of the null expansion $\theta_+$ is given by the stability operator $L_\Sigma$:
\begin{equation}
    \delta \theta_+[\phi] = L_\Sigma \phi = -\Delta_\Sigma \phi + 2\langle X, \nabla \phi \rangle + (Q + \Div_\Sigma X - |X|^2)\phi,
\end{equation}
where $X$ is the vector field dual to the one-form $k(\nu, \cdot)|_{T\Sigma}$ and $Q = \frac{1}{2}R_\Sigma - (\mu + J(\nu)) - \frac{1}{2}|\chi|^2$.

\subsection{The Variational Inequality (KKT Condition)}

\begin{theorem}[KKT Variational Inequality]\label{thm:KKT_Derivation}
If $\Sigma$ is a local maximizer of Area subject to $\theta_+ \le 0$, then for any variation $\phi$ that preserves the constraint to first order, the area must not increase.
The tangent cone of admissible variations at a point where $\theta_+ = 0$ is:
\begin{equation}
    \mathcal{T}_\Sigma = \{ \phi \in H^1(\Sigma) : L_\Sigma \phi \le 0 \text{ in the weak sense} \}.
\end{equation}
The first order optimality condition implies:
\begin{equation}
    \delta \mathcal{A}[\phi] \le 0 \quad \text{for all } \phi \in \mathcal{T}_\Sigma.
\end{equation}
Substituting $\delta \mathcal{A}[\phi] = \int_\Sigma H_\Sigma \phi \, dA$ and using the fact that $\theta_+ = H_\Sigma + \tr_\Sigma k = 0$ on the boundary of the trapped region (so $H_\Sigma = -\tr_\Sigma k$), we get:
\begin{equation}\label{eq:KKT_Condition_Derivation}
    -\int_\Sigma (\tr_\Sigma k) \phi \, dA \le 0 \quad \implies \quad \int_\Sigma (\tr_\Sigma k) \phi \, dA \ge 0 \quad \forall \phi \text{ s.t. } L_\Sigma \phi \le 0.
\end{equation}
This is the \textbf{Variational Inequality}. It is much stronger than the single eigenfunction condition $\int (\tr_\Sigma k) \psi_1 \ge 0$.
\end{theorem}

\begin{remark}[Fundamental Obstruction to Pointwise Jump]\label{rem:ConjectureCFundamental}
Proving the pointwise condition $\tr_\Sigma k \ge 0$ from the variational inequality \eqref{eq:KKT_Condition_Derivation} is likely impossible due to the maximum principle preventing the construction of sharply peaked supersolutions. The distributional condition is the natural limit of the variational principle.
\end{remark}

\subsection{Dual Formulation and Symmetrization}

By the generalized Lagrange Multiplier Theorem for convex optimization in Banach spaces (see, e.g., Luenberger \cite[Section 8.3, Theorem 1]{luenberger1969} or Zeidler \cite[Section 48.3]{zeidler1985}), the variational inequality is equivalent to the existence of a non-negative Lagrange multiplier measure $\mu \ge 0$ such that:
\begin{equation}
    -\tr_\Sigma k = L_\Sigma^* \mu,
\end{equation}
where $L_\Sigma^*$ is the formal adjoint of $L_\Sigma$ with respect to the $L^2$ inner product.

\textbf{Derivation of the Adjoint Operator:}
We compute the formal adjoint $L_\Sigma^*$ explicitly. Let $u, v \in C^\infty(\Sigma)$. Then:
\begin{align*}
    \langle u, L_\Sigma v \rangle_{L^2} &= \int_\Sigma u \left( -\Delta_\Sigma v + 2\langle X, \nabla v \rangle + (Q + \Div_\Sigma X - |X|^2)v \right) dA \\
    &= \int_\Sigma \left( -u \Delta_\Sigma v + 2u \langle X, \nabla v \rangle + u(Q + \Div_\Sigma X - |X|^2)v \right) dA.
\end{align*}
Integrating by parts (using that $\Sigma$ is closed):
\begin{itemize}
    \item Laplacian term: $\int_\Sigma -u \Delta_\Sigma v = \int_\Sigma \langle \nabla u, \nabla v \rangle = \int_\Sigma (-\Delta_\Sigma u) v$.
    \item Drift term: $\int_\Sigma 2u \langle X, \nabla v \rangle = \int_\Sigma 2 \langle uX, \nabla v \rangle = -\int_\Sigma 2 \Div_\Sigma(uX) v$.
\end{itemize}
Expanding the divergence term $\Div_\Sigma(uX) = \langle \nabla u, X \rangle + u \Div_\Sigma X$, the drift contribution becomes:
\[
    -\int_\Sigma 2 (\langle \nabla u, X \rangle + u \Div_\Sigma X) v = \int_\Sigma \left( -2\langle X, \nabla u \rangle - 2(\Div_\Sigma X)u \right) v.
\]
Combining all terms, we obtain the adjoint operator:
\begin{equation}
    L_\Sigma^* u = -\Delta_\Sigma u - 2\langle X, \nabla u \rangle + (Q - \Div_\Sigma X - |X|^2)u.
\end{equation}
Thus, the structural condition on the mean curvature jump is that it lies in the image of the positive cone under this adjoint operator.

\textbf{Derivation of the Symmetrization:}
To analyze the spectrum of this non-self-adjoint operator, we employ a symmetrization technique. We seek a function $\sigma$ such that the conjugated operator $\widetilde{L}_\Sigma := e^\sigma L_\Sigma e^{-\sigma}$ is self-adjoint.
Compute the action of the conjugated operator on a function $\psi$:
\[
    L_\Sigma(e^{-\sigma}\psi) = -\Delta(e^{-\sigma}\psi) + 2\langle X, \nabla(e^{-\sigma}\psi) \rangle + V e^{-\sigma}\psi,
\]
where $V = Q + \Div X - |X|^2$.
Using the identities $\nabla(e^{-\sigma}\psi) = e^{-\sigma}(\nabla\psi - \psi\nabla\sigma)$ and $\Delta(e^{-\sigma}\psi) = e^{-\sigma}(\Delta\psi - 2\langle\nabla\sigma, \nabla\psi\rangle + (|\nabla\sigma|^2 - \Delta\sigma)\psi)$, we find:
\begin{align*}
    e^\sigma L_\Sigma(e^{-\sigma}\psi) &= -(\Delta\psi - 2\langle\nabla\sigma, \nabla\psi\rangle + (|\nabla\sigma|^2 - \Delta\sigma)\psi) \\
    &\quad + 2\langle X, \nabla\psi - \psi\nabla\sigma \rangle + V\psi \\
    &= -\Delta\psi + 2\langle \nabla\sigma + X, \nabla\psi \rangle + (\Delta\sigma - |\nabla\sigma|^2 - 2\langle X, \nabla\sigma \rangle + V)\psi.
\end{align*}
To eliminate the non-self-adjoint drift term $2\langle \dots, \nabla\psi \rangle$, we require $X = -\nabla\sigma$. This is possible if $X$ is a gradient field (which is always true locally, and globally on $S^2$ if $X$ is closed). Assuming $X = -\nabla\sigma$, we have $\Div X = -\Delta\sigma$ and $\langle X, \nabla\sigma \rangle = -|X|^2$. Substituting these into the potential term:
\begin{align*}
    \widetilde{V} &= \Delta\sigma - |\nabla\sigma|^2 - 2(-|\nabla\sigma|^2) + (Q + \Div X - |X|^2) \\
    &= -\Div X - |X|^2 + 2|X|^2 + Q + \Div X - |X|^2 \\
    &= Q.
\end{align*}
Thus, if $X$ is a gradient, the symmetrized operator reduces to the Schr\"odinger operator $\widetilde{L}_\Sigma = -\Delta_\Sigma + Q$, which is manifestly self-adjoint. This allows us to use the spectral theory of self-adjoint operators to characterize the admissible jumps.

\subsection{Interface with AMO Test Functions}

We now state the precise interface between the KKT condition and the test functions required for the AMO monotonicity argument. This proposition clarifies exactly which class of functions satisfies the distributional favorable jump condition.

\begin{proposition}[Interface Lemma: KKT $\implies$ AMO Compatibility]\label{prop:KKT_AMO_Interface}
Let $\Sigma$ be a constrained area maximizer (so the KKT conditions hold). Let $w \in W^{1,2}(\Sigma)$ be a non-negative test function arising from the AMO level set method (specifically, $w = |\nabla u|^p|_\Sigma$).
If $w$ belongs to the cone of supersolutions for the stability operator:
\begin{equation}
    L_\Sigma w \le 0 \quad \text{in the weak sense},
\end{equation}
then the distributional mean curvature jump term satisfies the non-negativity condition:
\begin{equation}
    \langle [H]_{\bar{g}} \delta_\Sigma, w \rangle = \int_\Sigma (\tr_\Sigma k) w \, dA \ge 0.
\end{equation}
Consequently, the distributional scalar curvature of the smoothed Jang metric satisfies $R_{\hat{g}_\epsilon} \to \mu_{\text{dist}} \ge 0$ when tested against such functions.
\end{proposition}

\begin{proof}
The KKT condition for the constrained maximization problem (Theorem \ref{thm:KKT_Derivation}) implies the existence of a Lagrange multiplier measure $\mu \ge 0$ such that
\[ -\tr_\Sigma k = L_\Sigma^* \mu. \]
For any test function $w$, we have
\[ \int_\Sigma (-\tr_\Sigma k) w \, dA = \langle L_\Sigma^* \mu, w \rangle = \langle \mu, L_\Sigma w \rangle. \]
If $w$ is a supersolution ($L_\Sigma w \le 0$) and $\mu \ge 0$, then the pairing $\langle \mu, L_\Sigma w \rangle$ is non-positive (integral of a non-negative measure against a non-positive function). Thus:
\[ \int_\Sigma (-\tr_\Sigma k) w \, dA \le 0 \implies \int_\Sigma (\tr_\Sigma k) w \, dA \ge 0. \]
This confirms that the sign of the jump is favorable when integrated against any supersolution of the stability operator.
\end{proof}

\subsection{Spectral Characterization of Admissible Jumps}

Using the symmetrization, we can explicitly characterize the space of admissible mean curvature jumps.

\begin{proposition}[Interface Lemma: KKT to AMO]\label{prop:InterfaceLemma_Spectral}
Let $\Sigma$ be a local area maximizer subject to $\theta^+ \le 0$. Let $\mu \ge 0$ be the KKT multiplier satisfying $L_\Sigma^* \mu = -\tr_\Sigma k$.
Then for any "AMO test function" $w \in C^\infty(\Sigma)$ that lies in the admissible cone (specifically, any $w$ satisfying $L_\Sigma w \le 0$), we have the distributional favorable jump condition:
\begin{equation}
    \int_\Sigma (\tr_\Sigma k) w \, dA \ge 0.
\end{equation}
In particular, if the MOTS is strictly stable ($\lambda_1 > 0$), the condition holds for $w = -\psi_1$ (implying an integrated bound). If the MOTS is marginally stable ($\lambda_1 = 0$), it holds for $w = \psi_1$ (implying vanishing flux).
\end{proposition}

\begin{proof}
From the KKT condition, $-\tr_\Sigma k = L_\Sigma^* \mu$. Thus for any test function $w$:
\begin{align*}
    \int_\Sigma (\tr_\Sigma k) w \, dA &= \int_\Sigma (-L_\Sigma^* \mu) w \, dA \\
    &= -\int_\Sigma \mu (L_\Sigma w) \, dA.
\end{align*}
If $w$ is in the admissible cone, i.e., $L_\Sigma w \le 0$, then $-(L_\Sigma w) \ge 0$.
Since $\mu \ge 0$ (as a measure), the integral of a non-negative function against a non-negative measure is non-negative:
\[ \int_\Sigma \mu (-L_\Sigma w) \, dA \ge 0. \]
Therefore,
\[ \int_\Sigma (\tr_\Sigma k) w \, dA \ge 0. \]
This confirms that the KKT condition structurally guarantees the favorable sign for all test functions compatible with the variational constraint.
\end{proof}
Let $\{\psi_j\}_{j=1}^\infty$ be the orthonormal basis of eigenfunctions of the self-adjoint operator $\widetilde{L}_\Sigma$ with eigenvalues $\lambda_1 < \lambda_2 \le \dots$.
The KKT condition $-\tr_\Sigma k = L_\Sigma^* \mu$ can be rewritten using the conjugation relation $L_\Sigma^* = e^{-\sigma} \widetilde{L}_\Sigma e^\sigma$ (assuming $X = -\nabla \sigma$):
\begin{equation}
    -\tr_\Sigma k = e^{-\sigma} \widetilde{L}_\Sigma (e^\sigma \mu).
\end{equation}
Let $\tilde{\mu} = e^\sigma \mu$ be the weighted measure. Expanding $\tilde{\mu}$ in the eigenbasis (in the distributional sense):
\begin{equation}
    \tilde{\mu} = \sum_{j=1}^\infty c_j \psi_j, \quad \text{where } c_j = \langle \tilde{\mu}, \psi_j \rangle_{L^2}.
\end{equation}
Then the jump condition becomes:
\begin{equation}
    -\tr_\Sigma k = e^{-\sigma} \sum_{j=1}^\infty c_j \lambda_j \psi_j.
\end{equation}
The constraint $\mu \ge 0$ implies that $\tilde{\mu}$ is a positive measure. This imposes constraints on the coefficients $c_j$.
\begin{itemize}
    \item For the principal eigenfunction $\psi_1$ (which can be chosen positive), we have $c_1 = \int \psi_1 d\tilde{\mu} > 0$.
    \item If $\Sigma$ is strictly stable ($\lambda_1 > 0$), then all $\lambda_j > 0$, and the operator $\widetilde{L}_\Sigma$ is invertible. The jump is then "positive on average" in a spectral sense.
    \item If $\Sigma$ is marginally stable ($\lambda_1 = 0$), then the first term vanishes from the image, implying $\int (-\tr_\Sigma k) e^\sigma \psi_1 = 0$. This recovers the standard marginal stability condition.
\end{itemize}

\subsection{Why Pointwise Positivity Fails (The "Trap")}

It is tempting to try to prove $\tr_\Sigma k \ge 0$ pointwise by choosing a sequence of test functions $\phi_\epsilon$ approximating a Dirac delta $\delta_p$ at a point $p$.
However, the KKT condition requires $\phi_\epsilon$ to be a \textbf{supersolution} ($L_\Sigma \phi_\epsilon \le 0$).
For the standard Laplacian, a supersolution cannot have a strict local maximum (Maximum Principle). Thus, it is impossible to construct a smooth supersolution that is zero everywhere except for a positive peak at $p$.
The "best" one can do is the Green's function $G_p(x)$, which satisfies:
\begin{equation}
    L_\Sigma G_p = \delta_p.
\end{equation}
But this has the wrong sign! We need $L_\Sigma \phi \le 0$, so we would need $-G_p$, which is negative and singular.
Testing against $-G_p$ gives:
\begin{equation}
    \int (-\tr_\Sigma k) (-G_p) \ge 0 \implies \int (\tr_\Sigma k) G_p \ge 0.
\end{equation}
This proves that the \textbf{potential} generated by the jump is non-negative, not the jump itself.
\[
    (\widetilde{L}_\Sigma^{-1} (-\tr_\Sigma k))(p) \ge 0.
\]
This confirms that the pointwise positivity condition is likely false for general initial data, while the distributional condition (positivity of the potential) is the mathematically natural and correct statement.

\subsection{Distributional Compatibility with Smoothing}

The goal is to show that the distributional scalar curvature of the smoothed metric remains non-negative.
The scalar curvature distribution is:
\begin{equation}
    R_{\bar{g}} = R_{\text{bulk}} + 2[H]\delta_\Sigma.
\end{equation}
We need to show that for relevant test functions $u$ (e.g., $u = \phi^{-1}$ in the AMO argument):
\begin{equation}
    \langle R_{\bar{g}}, u \rangle = \int_{M \setminus \Sigma} R_{\text{bulk}} u \, dV + \int_\Sigma 2[H] u \, dA \ge 0.
\end{equation}

\textbf{Derivation of the Boundary Term:}
Using $[H] = \tr_\Sigma k$ and the KKT condition $-\tr_\Sigma k = L_\Sigma^* \mu$, we substitute into the boundary integral:
\begin{equation}
    \int_\Sigma 2[H] u \, dA = -2 \int_\Sigma (L_\Sigma^* \mu) u \, dA.
\end{equation}
By the definition of the adjoint operator, we transfer the operator to $u$:
\begin{equation}
    -2 \int_\Sigma (L_\Sigma^* \mu) u \, dA = -2 \int_\Sigma \mu (L_\Sigma u) \, dA.
\end{equation}
Since $\mu$ is a non-negative measure (guaranteed by the KKT condition), the sign of this term is determined entirely by the sign of $L_\Sigma u$.
\begin{itemize}
    \item If $u$ is a \textbf{supersolution} ($L_\Sigma u \le 0$), then the product $\mu (L_\Sigma u)$ is non-positive (measure $\times$ non-positive function).
    \item The factor $-2$ flips the sign, making the total contribution non-negative:
    \[
        -2 \int_\Sigma \underbrace{\mu}_{\ge 0} \underbrace{(L_\Sigma u)}_{\le 0} \ge 0.
    \]
\end{itemize}
This derivation proves that the "Distributional Favorable Jump" condition is exactly equivalent to requiring non-negativity against supersolutions. Since the test functions in the AMO proof are constructed from level sets of Green's functions (which are supersolutions), the KKT condition structurally guarantees the validity of the inequality.

The "Minimal Distributional Upgrade" (formerly Conjecture 34.2) is thus resolved by Theorem D: for the outermost MOTS, the KKT condition provides the strongest possible structural guarantee, ensuring the inequality holds without additional assumptions.

\subsection{Proof of Theorem D: Distributional Favorable Jump}
\label{subsec:ProofTheoremD}

We explicitly verify that the weight function appearing in the distributional Bochner identity satisfies the supersolution condition required by the KKT argument. This constitutes the proof of Theorem D.

\textbf{The Weight Function:}
In the AMO monotonicity proof (Theorem \ref{thm:AMOMonotonicity}), the scalar curvature term arises from the weighted Bochner identity:
\begin{equation}
    \int_M |\nabla u|^{p-2} \Ric(\nabla u, \nabla u) \varphi \, dV.
\end{equation}
Near the boundary $\Sigma = \{u=0\}$, the gradient $\nabla u$ is parallel to the normal $\nu$. Thus, the effective weight function on $\Sigma$ is:
\begin{equation}
    w := |\nabla u|^p \big|_\Sigma.
\end{equation}
We must check the sign of $\langle [H]\delta_\Sigma, w \rangle$. By the KKT condition, this is non-negative if $w$ is a supersolution of the stability operator:
\begin{equation}
    L_\Sigma w \le 0.
\end{equation}

\textbf{The Calculation:}
For a $p$-harmonic function $u$, the gradient modulus $|\nabla u|$ satisfies a refined Kato inequality. Specifically, on a level set $\Sigma$, we have:
\begin{equation}
    \Delta_\Sigma |\nabla u| \ge |\nabla u| V + \frac{|\nabla |\nabla u||^2}{|\nabla u|},
\end{equation}
where $V = |A|^2 + \Ric(\nu, \nu)$.
The stability operator is $L_\Sigma = -\Delta_\Sigma - V$.
Applying this to $|\nabla u|$:
\begin{align*}
    L_\Sigma |\nabla u| &= -\Delta_\Sigma |\nabla u| - V |\nabla u| \\
    &\le - |\nabla u| V - V |\nabla u| = - 2 V |\nabla u|.
\end{align*}
Under the Dominant Energy Condition, $V \ge 0$.
Now consider the actual weight $w = |\nabla u|^p$.
\begin{align*}
    L_\Sigma (|\nabla u|^p) &= -\Delta_\Sigma (|\nabla u|^p) - V |\nabla u|^p \\
    &= -p |\nabla u|^{p-1} \Delta_\Sigma |\nabla u| - p(p-1) |\nabla u|^{p-2} |\nabla |\nabla u||^2 - V |\nabla u|^p.
\end{align*}
Using $-\Delta_\Sigma |\nabla u| \le -V |\nabla u|$ (from the Kato inequality step):
\begin{align*}
    L_\Sigma (|\nabla u|^p) &\le p |\nabla u|^{p-1} (-V |\nabla u|) - p(p-1) |\nabla u|^{p-2} |\nabla |\nabla u||^2 - V |\nabla u|^p \\
    &= -(p+1) V |\nabla u|^p - p(p-1) |\nabla u|^{p-2} |\nabla |\nabla u||^2.
\end{align*}
Since $V \ge 0$ and $p \ge 1$, both terms are non-positive. Thus, $L_\Sigma w \le 0$ holds.

The test function $w = |\nabla u|^p$ inherits the supersolution property from the Bochner identity. This closes the logical loop: the $p$-Laplacian that generates the monotonicity also generates the test weights to unlock the KKT positivity.

\subsection{Verification of the supersolution condition}
\label{subsec:Verification_KKT}

We verify that $w = |\nabla u|^p|_\Sigma$ satisfies $L_\Sigma w \le 0$. The boundary term in the AMO monotonicity formula is $\int_\Sigma [H] |\nabla u|^p \, dA$, so $w = |\nabla u|^p$ is the relevant test function. The $p$-harmonic function $u$ satisfies the refined Kato inequality
\[
\Delta_\Sigma |\nabla u| \ge |\nabla u| (|A|^2 + \Ric(\nu, \nu)) + \frac{|\nabla |\nabla u||^2}{|\nabla u|}.
\]
With $L_\Sigma = -\Delta_\Sigma - (|A|^2 + \Ric(\nu, \nu))$ and $V = |A|^2 + \Ric(\nu,\nu)$, we obtain $L_\Sigma |\nabla u| \le -2V|\nabla u| \le 0$. For $w = |\nabla u|^p$,
\[
L_\Sigma w \le -(p+1) V w - p(p-1) |\nabla u|^{p-2} |\nabla |\nabla u||^2 \le 0,
\]
confirming the supersolution property. The KKT condition and AMO method are thus compatible.

\subsection{Distributional compatibility}

A potential objection to the KKT upgrade is whether the existence of a measure $\mu$ is sufficient to control the scalar curvature in the distributional sense required for the smoothing argument. We address this explicitly:

\begin{remark}[Distributional vs.\ Pointwise Positivity]
The standard Jang reduction requires $\tr_\Sigma k \ge 0$ pointwise to ensure the scalar curvature of the Jang metric is non-negative. In our generalized setting, the scalar curvature appears as a distribution $R_{\bar{g}} = R_{\bar{g}}^{reg} + 2[H]_{\bar{g}} \delta_\Sigma$.
The KKT condition provides a measure $\mu \ge 0$ such that $-\tr_\Sigma k = L_\Sigma^* \mu$.
Crucially, the smoothing procedure (Miao's method) does not require pointwise non-negativity of the jump. It requires that the jump is "distributionally non-negative" in the sense that it can be approximated by smooth metrics with non-negative scalar curvature.
The existence of $\mu \ge 0$ ensures exactly this: the negative part of the jump is in the image of the adjoint stability operator acting on a positive measure. This structure allows us to deform the metric in a neighborhood of $\Sigma$ to absorb the negative contribution into the bulk scalar curvature, preserving the global non-negativity condition in the limit.
\end{remark}
